\section{Time weighting for the flow loss}
\label{app:flow-weighting}

The flow-matching term can be reweighted in time by introducing a scalar
$\alpha(t)$:
\begin{equation}
\mathcal{L}_{\text{flow},\alpha}(\theta)
=
\mathbb{E}_{(x_0,x_1)\sim\tilde{\gamma},\,t\sim\Unif([0,1])}\left[
\alpha(t)\,\|-\nabla_{\hat{x}} V_\theta(\hat{x}_t)-v\|_2^2
\right].
\end{equation}
A simple choice is a piecewise linear schedule
\begin{equation}
\alpha_{\tau^\ast}(t)=
\begin{cases}
1,&0\le t\le\tau^\ast,\\[6pt]
\dfrac{1-t}{1-\tau^\ast},&\tau^\ast<t\le1,
\end{cases}
\end{equation}
so that $\alpha_{\tau^\ast}(1)=0$. This downweights flow supervision for samples
approaching the data manifold and can reduce interference with $\mathcal{L}_{\text{ER}}$.
We do not apply this weighting in our experiments ($\alpha(t)\equiv 1$), because
full flow supervision was empirically more stable for $\mathcal{L}_{\text{ER}}$.

Time-decaying weights are common in continuous settings: Energy Matching and
Equilibrium Matching relax the flow term to avoid over-constraining the
equilibrium regime. In particular, Equilibrium Matching uses $\alpha(t)\to 0$
as $t\to 1$ because a vanishing energy gradient at equilibrium is one of the
score-matching conditions for energy-based training~\citep{wang2025equilibrium}.
Such schedules may be beneficial in other domains where late-time supervision
conflicts with equilibrium refinement.
